\section{Experiments}
\label{sec::appendix}

\subsection{Finite-state diagnostic experiments}

\paragraph{Baird hub-and-spokes.}
The Baird experiment uses a hub state and six spoke states
\citep{baird1995residual}. The solid action moves deterministically to the hub;
the dashed action moves uniformly to one of the spokes. The target policy always
takes the solid action. The behavior policy takes the solid action with
probability \(\kappa\), so decreasing \(\kappa\) reduces stationary overlap with
the target policy. We use the standard seven-dimensional linear feature map,
zero rewards, ridge-regularized least squares at each FQE step, and exact
stationary ratios. Since \(Q^\pi\equiv0\), the experiment isolates whether the
projected Bellman iteration contracts in the intended norm. Reported curves are
averaged over 100 independent datasets.

\paragraph{Sparse Garnet MDPs.}
The Garnet experiment uses random MDPs with \(100\) states, \(4\) actions, and
branching factor \(5\). For each state-action pair, five next states are sampled
without replacement and assigned Dirichlet transition probabilities. Rewards
are Gaussian. The target policy is sampled randomly; the behavior policy is
\(\varepsilon\)-greedy with \(\varepsilon=0.1\). Data are generated with a reset
scheme that deliberately breaks stationarity. The linear feature class has
dimension \(5\) and contains \(Q^\pi\), but is generically not Bellman complete.
Thus any instability is due to the fitted projection geometry, not ordinary
realizability failure. The main figure uses exact stationary ratios; the same
qualitative stabilization was also observed with tabular DualDICE-style ratio
estimates.

\subsection{Controlled linear-Gaussian benchmark}

The continuous benchmark uses \(S_t\in\mathbb R^2\) and \(A_t\in\mathbb R\).
The dynamics are linear Gaussian,
\[
S_{t+1}=B S_t + C A_t+\epsilon_t,
\qquad
\epsilon_t\sim N(0,\sigma^2 I),
\]
with quadratic reward
\[
r(s,a)=-\|s-s^\star\|^2-\lambda a^2 .
\]
The target and behavior policies are Gaussian linear policies,
\[
A_t\mid S_t=s \sim N(K_\pi s,\sigma_\pi^2),
\qquad
A_t\mid S_t=s \sim N(K_b s,\sigma_b^2),
\]
where \(K_b=K_\pi+\Delta v\). The scalar \(\Delta\) is the shift shown in the
plots. Low, moderate, and severe shifts correspond to \(\Delta=0,1.1,1.35\).

Because the system is linear Gaussian with quadratic rewards, \(Q^\pi\),
\(V^\pi\), and the discounted value \(\psi_\pi\) are computed analytically by a
quadratic fixed-point recursion. Discounted state-action occupancies are
represented as Gaussian mixtures over time. For \(\gamma_{\rm ratio}=1\), the
reference distribution is the stationary Gaussian state-action distribution
under the corresponding policy. Oracle weights are evaluated as stabilized
mixture-density or stationary Gaussian density ratios.

Offline data are sampled from the behavior reference distribution matched to
\(\gamma_{\rm ratio}\). For \(\gamma_{\rm ratio}<1\), this means drawing a
geometric time \(T\) and recording a transition from the behavior trajectory at
time \(T\). For \(\gamma_{\rm ratio}=1\), data are sampled from the behavior
stationary distribution. The value estimand is always the discounted
\(\gamma=0.95\) policy value; \(\gamma_{\rm ratio}\) only changes the regression
weighting target. This distinction is important: stationary weighting can be
excellent for stationary-reference \(Q\) error while not necessarily optimizing
the discounted initial-state value.

FQE is run with three feature regimes. The main text uses the misspecified
affine class \([1,s_1,s_2,a]\). A full quadratic class serves as a
well-specified check, and a diagonal quadratic class is used as an additional
misspecification stress test. Linear FQE uses analytic target-action
expectations, so there is no Monte Carlo noise in the Bellman target. Neural FQE
uses a small \(Q(s,a)\) network and Gauss-Hermite quadrature for the target
action expectation.

The ratio estimators use the discounted or stationary flow moment equations.
For \(\gamma_{\rm ratio}<1\), the moment condition includes the
\((1-\gamma_{\rm ratio})\nu_0\pi\) right-hand side. For
\(\gamma_{\rm ratio}=1\), the stationary moment has zero right-hand side plus a
mean-one normalization. We report raw oracle weights, clipped oracle weights,
linear estimated weights, clipped estimated weights, and neural/RKHS estimated
weights. Clipping is fixed in advance by an ESS target and maximum-weight cap;
it is not tuned separately by shift or estimator.

The primary metrics are target-reference \(Q\) MSE, behavior-reference \(Q\)
MSE, and discounted policy-value MSE. Diagnostics include ESS fraction, weight
quantiles and maxima, estimated-vs-oracle log-ratio RMSE, correlation with
oracle weights, feature calibration error, unstable-run fraction, and Monte
Carlo standard errors across seeds. The final paper run uses \(100\) seeds for
the main linear estimators and \(20\) seeds for neural-ratio and neural-FQE
variants.

\begin{table}[t]
\centering
\small
\setlength{\tabcolsep}{3.5pt}
\caption{Controlled linear-Gaussian benchmark with misspecified affine FQE, $n=4000$, value discount $\gamma=0.95$, and discounted weighting target $\gamma_{\rm ratio}=0.95$. Entries report target-reference $Q$ MSE / policy-value MSE / ESS fraction / weight $q_{.99}$. Unweighted FQE has ESS fraction and $q_{.99}$ equal to one by construction.}
\label{tab:controlled-benchmark}
\resizebox{\linewidth}{!}{%
\begin{tabular}{lrrrr}
\toprule
Shift & Standard & Oracle clipped & Estimated clipped & Neural-ratio clipped \\
\midrule
Low $(0)$ &
$3.27/0.17/1.00/1.00$ &
$3.27/0.17/1.00/1.00$ &
$3.21/0.11/0.995/1.18$ &
$2.94/0.26/0.997/1.13$ \\
Moderate $(1.1)$ &
$505/132/1.00/1.00$ &
$71.0/7.69/0.400/7.75$ &
$16.2/66.4/0.400/4.96$ &
$22.7/115/0.400/5.52$ \\
Severe $(1.35)$ &
$1.34{\times}10^4/9.87{\times}10^3/1.00/1.00$ &
$1.37{\times}10^3/1.04{\times}10^3/0.142/21.8$ &
$1.83{\times}10^3/2.46{\times}10^3/0.400/7.30$ &
$1.35{\times}10^3/884/0.400/7.00$ \\
\bottomrule
\end{tabular}
}
\end{table}


The controlled benchmark supports a narrower but more useful claim than
uniform dominance. Under low shift, all methods are similar. Under moderate
shift with misspecified affine FQE, weighting sharply reduces target-reference
\(Q\) error and improves value accuracy. Under severe shift, the same weighting
principle still helps target-reference error, but ESS and upper weight quantiles
show that the problem is close to an overlap boundary. Neural FQE reduces the
need for weighting because the function class is more flexible; nevertheless,
near-stationary discounted weighting with \(\gamma_{\rm ratio}=0.99\) improves
both target-reference \(Q\) error and value MSE in the moderate and severe
shift settings. We view this as evidence that weighting is most useful when
approximation error is distribution-sensitive, and as a reminder that ratio
estimation and overlap remain practical bottlenecks.
